\subsection{Условия}

Ветвление записывается словами \texttt{if} \dots\ \texttt{then} \dots\ \texttt{elseif} \dots\ \texttt{then} \dots\ \texttt{else} \dots\ \texttt{end}. Ветка \verb|elseif| (слитно!) может повторяться, \verb|else| --- не обязательна, \verb|end| --- обязателен всегда.

\begin{codefile}{lua/04_language/lua_conditions_if.lua}
\end{codefile}

\paragraph{Что происходит.} Функция \verb|classify| проверяет условия по порядку и возвращает первое подходящее слово. Для \verb|0| первая проверка \verb|x > 0| ложна, вторая \verb|x == 0| истинна --- получаем \verb|"zero"|.

\subsubsection*{Охранное условие}

Обращение к полю несуществующей таблицы --- ошибка «attempt to index a nil value». Защищаются от неё проверкой перед обращением:

\begin{codefile}{lua/04_language/lua_conditions_check.lua}
\end{codefile}

Условие \verb|t and t.value| читается так: «\verb|t| существует, и у неё есть поле \verb|value|». Если \verb|t| равно \verb|nil|, до \verb|t.value| дело не доходит --- спасает короткое замыкание \verb|and|.

\begin{hintbox}[Таблица вместо каскада if]
Когда вариантов много и каждому соответствует действие, удобнее таблица: \verb|colors[state]| или даже таблица функций \verb|actions[state]()|. Так в главе~\ref{sec:leds} устроен светофор --- самый маленький конечный автомат.
\end{hintbox}

\begin{taskbox}[Две ветки, три проверки]
\begin{enumerate}
  \item Флаг \verb|ready = false|. Напечатайте \verb|"красный"|, если флаг ложен, и \verb|"зелёный"|, если истинен. Запустите дважды, поменяв только значение флага.
  \item Порог высоты --- 2~м. Напечатайте \verb|"выше"|, только если \verb|height| \textbf{строго} больше порога, иначе \verb|"не выше"|. Проверьте на 1, 2 и 3: ниже границы, на границе и выше.
  \item Тройное ветвление по температуре: \verb|temp < 0| $\to$ \verb|"cold"|, от 0 до 20 (не включая 20) $\to$ \verb|"cool"|, иначе \verb|"warm"|. Проверьте на $-5$, 0, 20.
  \item Выберите цвет: \verb|{1,0,0}| для \verb|state == "RED"|, \verb|{1,1,0}| для \verb|"YELLOW"|, иначе \verb|{0,1,0}|; напечатайте три составляющие через \verb|table.unpack|.
  \item Защитите деление от нуля: если \verb|b ~= 0|, напечатайте \verb|a / b|, иначе --- сообщение «делить на ноль нельзя».
\end{enumerate}
\expected{1) \texttt{красный}, после замены флага --- \texttt{зелёный}; 2) \texttt{не выше}, \texttt{не выше}, \texttt{выше} --- граница 2 к «выше» не относится; 3) \texttt{cold}, \texttt{cool}, \texttt{warm}; 4) для \texttt{"YELLOW"} --- \texttt{1 1 0}; 5) при \texttt{b = 0} --- сообщение, а не \texttt{inf}.}
\mission{Lua}{22--23}
\end{taskbox}
